\documentclass[12pt,a4paper]{article}
\usepackage[utf8]{inputenc}
\usepackage[T5]{fontenc}          % T5 = Vietnamese encoding (KHÔNG dùng T1)
\usepackage[vietnamese]{babel}
\usepackage{amsmath, amssymb, amsfonts}
\usepackage{soul}

\begin{document}

\ul{Buổi 05 -- Ngày 16-08-2026 -- môn Đại số tuyến tính -- lớp
MA003.F32.LT.CNTT}

\textbf{Chương 3: KHÔNG GIAN VÉC TƠ}

\textbf{1/ \ul{KHÔNG GIAN VÉC TƠ
}}$\mathbb{R}^n$ \textbf{\ul{VÀ
KHÔNG GIAN VÉC TƠ TỔNG QUÁT}}

Ta định nghĩa \[ \mathbb{R}^n = \mathbb{R} \times \mathbb{R} \times \dots \times \mathbb{R} = \{(x_1; x_2; \dots; x_n) \mid x_1, x_2, \dots, x_n \in \mathbb{R}\} \]

$n$ lần

\ul{Ví dụ}:

\begin{align*}
 \mathbb{R}^1 &= \mathbb{R} \\ \mathbb{R}^2 &= \mathbb{R} \times \mathbb{R} = \{\alpha = (a;b) \mid a,b \in \mathbb{R}\} \\ \mathbb{R}^3 &= \mathbb{R} \times \mathbb{R} \times \mathbb{R} = \{\alpha = (a;b;c) \mid a,b,c \in \mathbb{R}\} 
\end{align*}

Trên $\mathbb{R}^n$ ta định
nghĩa 2 phép toán:

a/ \emph{\ul{Phép cộng véc tơ}}:

$\forall \alpha = (x_1; x_2; \dots; x_n) \in \mathbb{R}^n$, và
$\forall \beta = (y_1; y_2; \dots; y_n) \in \mathbb{R}^n$, ta có

$\alpha + \beta = (x_1 + y_1; x_2 + y_2; \dots; x_n + y_n) \in \mathbb{R}^n$

Ta cần nhớ quy tắc:

b/ \emph{\ul{Phép nhân một con số thực với véc tơ}}

Với $\forall \alpha = (x_1; x_2; \dots; x_n) \in \mathbb{R}^n$ và số thực
$c \in \mathbb{R}$, ta có:

$c.\alpha = (cx_1; cx_2; \dots; cx_n) \in \mathbb{R}^n$

\emph{\ul{Lưu ý}}: véc tơ
$c.\alpha$ là véc tơ có
cùng phương (cùng giá đỡ -- support) với véc tơ
$\alpha$.

và $c.\alpha$ có cùng
chiều với $\alpha$ nếu
$c > 0$ và
$c.\alpha$ ngược chiều với
$\alpha$ nếu
$c < 0$.

Khi $c = 0$ thì véc tơ
$c.\alpha$ là véc tơ zero.

Đồng thời, véc tơ
$c.\alpha$ có độ dài gấp
$|c|$ lần độ dài của
véc tơ $\alpha$.

Ta gọi cấu trúc đại số:
$(\mathbb{R}^n, +, \bullet)$ là không gian
tuyến tính (không gian véc tơ)
$n$ chiều trên
trường số thực $\mathbb{R}$.

Tổng quát, ta ký hiệu
$(V, +, \bullet)$ là không gian
tuyến tính (không gian véc tơ) tổng quát trên trường số \emph{F}, với
phép cộng và phép nhân được định nghĩa tương tự như trên
$\mathbb{R}^n$ (\emph{V} =
Vector space).

\ul{Ví dụ}: Trên
$\mathbb{R}^6$cho
\[ \begin{cases} \alpha = (-2; 3; 5; -8; -4; 9) \\ \beta = (7; -1; 2; 4; -6; 3) \end{cases} \]. Tìm
\[ \begin{cases} 235\alpha - 1479\beta = ? \\ 3572\beta - 5383\alpha = ? \end{cases} \]

\textbf{2/ \ul{KHÔNG GIAN VÉC TƠ CON (SUB-VECTOR SPACE)}}

Cho không gian véc tơ
$(V, +, \bullet)$ là không gian
véc tơ tổng quát trên trường số \emph{F},

Và cho tập hợp $W \neq \emptyset$
trên $V$.

$W$

Lúc này do $W \subset V$ nên
các véc tơ trên $W$
được thừa hưởng phép cộng (+) và nhân
()
có sẵn trên $V$.

Ta nói $(W, +, \bullet)$ là không
gian tuyến tính (không gian véc tơ) con của
$(V, +, \bullet)$

và ký hiệu là $(W, +, \bullet) \le (V, +, \bullet)$,
và viết gọn là $W \le V$,
nếu ta chứng minh được 2 tính chất sau:

1/ $\forall \alpha \in W, \forall \beta \in W$, ta có
$(\alpha + \beta) \in W$.

2/ $\forall \alpha \in W, \forall c \in F$, ta chứng tỏ
$(c.\alpha) \in W$.

Ngược lại, nếu 1 trong 2 tính chất này bị vi phạm thì
$W$ không phải là
không gian véc tơ con của
$V$, và ta kí hiệu
là $W \le V$

\ul{Ví dụ mẫu 1}: Trên
$V = \mathbb{R}^3$ cho tập hợp
\[ W = \{(x; y; z) \in \mathbb{R}^3 \mid 2x - y + 3z = 0\} \].

Hỏi $W \le V$ không? Vì
sao?

\ul{Ví dụ mẫu 2}: Trên
$V = \mathbb{R}^3$ cho tập hợp
\[ W = \{(x; y; z) \in \mathbb{R}^3 \mid x^2 - 5y + 8z = 1\} \].

Hỏi $W \le V$ không? Vì
sao?

\ul{Ví dụ mẫu 3}: Trên
$V = M_2(\mathbb{R})$ là không gian
các ma trận vuông, thực, cấp 2, cho tập hợp

\[ W = \left\{ A = \begin{pmatrix} a & -b \\ 2b & 4a \end{pmatrix} \mid a, b \in \mathbb{R} \right\} \]

Hỏi $W \le V$ không? Vì
sao?

\ul{Ví dụ mẫu 4}: Trên
$V = M_2(\mathbb{R})$ là không gian
các ma trận vuông, thực, cấp 2, cho tập hợp

\[ W = \left\{ A = \begin{pmatrix} a^2 & b \\ ab & -2a \end{pmatrix} \mid a, b \in \mathbb{R} \right\} \]

Hỏi $W \le V$ không? Vì
sao?

\ul{Giải}:

\ul{Ví dụ mẫu 1}:
\[ W = \{(x; y; z) \in \mathbb{R}^3 \mid 2x - y + 3z = 0\} \]

Ta chứng tỏ $W \le V$ như
sau:

Ta có $\gamma = (1; -1; -1) \in W$ do 2.1 --
(--1) + 3.(--1) = 0

$\Rightarrow W \neq \emptyset$ (*)

Mặt khác:

$\forall \alpha = (x_1; y_1; z_1) \in W$ ta có
$2x_1 - y_1 + 3z_1 = 0$ (1)

$\forall \beta = (x_2; y_2; z_2) \in W$ ta có
$2x_2 - y_2 + 3z_2 = 0$ (2)

Cộng (1), (2) theo vế ta được:

$(2x_1 - y_1 + 3z_1) + (2x_2 - y_2 + 3z_2) = 0$

$\Leftrightarrow 2(x_1 + x_2) - (y_1 + y_2) + 3(z_1 + z_2) = 0$

Mà $\alpha + \beta = (x_1 + x_2; y_1 + y_2; z_1 + z_2) \Rightarrow (\alpha + \beta) \in W$ (**)

Ngoài ra, $\forall c \in \mathbb{R}$ ta
nhân 2 vế pt (1) cho
$C$ thì được:

$c(2x_1 - y_1 + 3z_1) = 0$

Ta nhân $C$ với
$\forall \alpha = (x_1; y_1; z_1) \in W$ thì có
$c\alpha = c(x_1; y_1; z_1) = (cx_1; cy_1; cz_1)$

$\Rightarrow (c\alpha) \in W$ (***)

Từ (*), (**), (***)
$\Rightarrow W \le V$.

\ul{Ví dụ mẫu 2}:
\[ W = \{(x; y; z) \in \mathbb{R}^3 \mid x^2 - 5y + 8z = 1\} \]

Ta chứng tỏ $W \le V$ như
sau:

\ul{Cách 1}:

Ta chọn $\alpha = (4; 3; 0) \in W$ do
$4^2 - 5.3 + 8.0 = 1$

Và chọn $\beta = (-4; 3; 0) \in W$do
$(-4)^2 - 5.3 + 8.0 = 1$

Mà $\alpha + \beta = (0; 6; 0)$ có
$0^2 - 5.6 + 8.0 = -30 \neq 1$ nên
$(\alpha + \beta) \notin W$

Cho nên do $W \le V$.

\ul{Cách 2}:

Ta chọn $\alpha = (3; 0; -1) \in W$ do
$3^2 - 5.0 + 8.(-1) = 1$

Và chọn $c = 2 \Rightarrow c\alpha = (6; 0; -2)$ có
$6^2 - 5.0 + 8.(-2) = 36 - 16 = 20 \neq 1$

nên $(c\alpha) \notin W$ . Cho nên
ta có $W \not\le V$.

\ul{Ví dụ mẫu 3}: Trên
$V = M_2(\mathbb{R})$ là không gian
các ma trận vuông, thực, cấp 2, cho tập hợp

\[ W = \left\{ A = \begin{pmatrix} a & -b \\ 2b & 4a \end{pmatrix} \mid a, b \in \mathbb{R} \right\} \]

Ta chứng tỏ $W \not\le V$ như
sau:

Ta chọn $a=1, b=3$, ứng
với $a=a_1, b=b_1$

$\Rightarrow W \neq \emptyset$ (*)

Mặt khác:

$a=a_2, b=b_2$ ứng với
$\forall \alpha = A_1 = \begin{pmatrix} a_1 & -b_1 \\ 2b_1 & 4a_1 \end{pmatrix} \in W$

$\forall \beta = A_2 = \begin{pmatrix} a_2 & -b_2 \\ 2b_2 & 4a_2 \end{pmatrix} \in W$ ứng với
$a=a_2, b=b_2$

Suy ra \[ \alpha + \beta = A_1 + A_2 = \begin{pmatrix} a_1+a_2 & -b_1-b_2 \\ 2b_1+2b_2 & 4a_1+4a_2 \end{pmatrix} = \begin{pmatrix} a_1+a_2 & -(b_1+b_2) \\ 2(b_1+b_2) & 4(a_1+a_2) \end{pmatrix} \] ứng với

$a=a_1+a_2, b=b_1+b_2$

Cho nên $(\alpha + \beta) \in W$ (**)

Ngoài ra $\forall c \in \mathbb{R}$, ta có

\begin{quote}
\[ c\alpha = cA_1 = c \begin{pmatrix} a_1 & -b_1 \\ 2b_1 & 4a_1 \end{pmatrix} = \begin{pmatrix} ca_1 & -cb_1 \\ c(2b_1) & c(4a_1) \end{pmatrix} = \begin{pmatrix} ca_1 & -(cb_1) \\ 2(cb_1) & 4(ca_1) \end{pmatrix} \]
\end{quote}

ứng với

$a = ca_1, b = cb_1$

Cho nên $(c\alpha) \in W$ (***)

Từ (*), (**), (***)
$\Rightarrow W \le V$.

\ul{Ví dụ mẫu 4}: Trên
$V = M_2(\mathbb{R})$ là không gian
các ma trận vuông, thực, cấp 2, cho tập hợp

\[ W = \left\{ A = \begin{pmatrix} a^2 & b \\ ab & -2a \end{pmatrix} \mid a, b \in \mathbb{R} \right\} \]

Ta chứng tỏ $W \le V$ như
sau:

Ta chọn $a=2, b=1$ suy ra
\[ \alpha = A_1 = \begin{pmatrix} 2^2 & 1 \\ 2.1 & -2.2 \end{pmatrix} = \begin{pmatrix} 4 & 1 \\ 2 & -4 \end{pmatrix} \in W \]

Và chọn $a=-2, b=1$ suy ra
\[ \beta = A_2 = \begin{pmatrix} (-2)^2 & 1 \\ -2.1 & -2.(-2) \end{pmatrix} = \begin{pmatrix} 4 & 1 \\ -2 & 4 \end{pmatrix} \in W \]

Mà \[ \alpha + \beta = A_1 + A_2 = \begin{pmatrix} 8 & 2 \\ 0 & 0 \end{pmatrix} \notin W \] do không còn
thỏa điều kiện của \emph{W}.

(do \[ \begin{cases} a^2 = 8 \\ -2a = 0 \end{cases} \Leftrightarrow \begin{cases} a = \pm 2\sqrt{2} \\ a = 0 \end{cases} \] vô lí)

Cho nên $W \not\le V$.

\ul{Bài tập}:

Kiểm chứng \emph{W} có phải là không gian con của \emph{V} hay không? Vì
sao?

a/ $V = \mathbb{R}^3, W = \{(x; y; z) \mid 3x + 5z = 0\}$.

b/ $V = \mathbb{R}^3, W = \{(x; y; z) \mid y = 0\}$.

c/ $V = \mathbb{R}^3, W = \{(x; y; z) \mid 4x = 2y - z\}$.

d/ $V = \mathbb{R}^3, W = \{(x; y; z) \mid 6x - 5yz = 0\}$.

e/ $V = \mathbb{R}^3, W = \{(x; y; z) \mid xyz = 0\}$.

f/ $V = \mathbb{R}^3, W = \{(x; y; z) \mid 7x - z = 2\}$.

g/ $V = \mathbb{R}^3, W = \{(x; y; z) \mid 2x - 5z^2 = 0\}$.

h/ $V = \mathbb{R}^4, W = \{(x; y; z; t) \mid 2x - 4y + z - 6t = 0\}$.

i/ \[ V = \mathbb{R}^4, W = \begin{cases} x + 5z = 0 \\ 4y - z = 0 \end{cases} \].

j/ \[ V = \mathbb{R}^4, W = \left\{ (x;y;z;t) \mid \begin{cases} x + y - 2z + 3t = 0 \\ x = 4z - t \end{cases} \right\} \].

k/ $V = M_2(\mathbb{R})$ là không
gian các ma trận vuông, thực, cấp 2, cho tập hợp

\[ W = \left\{ A = \begin{pmatrix} 3a & 5b \\ -b & a+b \end{pmatrix} \mid a,b \in \mathbb{R} \right\} \].

l/ $V = M_2(\mathbb{R})$ là không
gian các ma trận vuông, thực, cấp 2, cho tập hợp

\[ W = \left\{ A = \begin{pmatrix} ab & 2b \\ b^3 & a \end{pmatrix} \mid a,b \in \mathbb{R} \right\} \].

\textbf{3/ \ul{TỔ HỢP TUYẾN TÍNH (LINEAR COMBINATION)}}

Cho $(V, +, \bullet)$ là không
gian tuyến tính tổng quát và cho các véc tơ
$\alpha_1, \alpha_2, \dots, \alpha_m$ trên \emph{V}.

Ta chọn bộ số $c_1, c_2, \dots, c_m$
trên \emph{F} và lập véc tơ

$\alpha = c_1\alpha_1 + c_2\alpha_2 + \dots + c_m\alpha_m$.

Thì ta gọi $\alpha$ là
một tổ hợp tuyến tính (linear combination) của
$\alpha_1, \alpha_2, \dots, \alpha_m$.

Do các số $c_1, c_2, \dots, c_m$ được
chọn tùy ý nên ta có thể lập được vô số các tổ hợp tuyến tính (THTT) của
$\alpha_1, \alpha_2, \dots, \alpha_m$.

Gọi $S = \{\alpha_1, \alpha_2, \dots, \alpha_m\}$ là tập hợp
trên \emph{V}.

Thì véc tơ $\alpha = c_1\alpha_1 + c_2\alpha_2 + \dots + c_m\alpha_m$
được gọi là một tổ hợp tuyến tính có từ \emph{S}.

\ul{Ví dụ}:

Trên $V = \mathbb{R}^5$, ta cho
\[ S = \{\alpha_1 = (-5; 2; -4; 2; 6), \alpha_2 = (7; -1; 8; 6; 3), \alpha_3 = (2; -1; -2; 4; 9)\} \].

a/ Chọn $c_1 = -2, c_2 = 3, c_3 = 1$, ta có
một tổ hợp tuyến tính có từ \emph{S} là:

$\alpha = c_1\alpha_1 + c_2\alpha_2 + c_3\alpha_3$

\begin{align*}
 &= -2(-5, 2, -4, 2, 6) + 3(7, -1, 8, 6, 3) + 1(2, -1, -2, 4, 9) \\ &= (10, -4, 8, -4, -12) + (21, -3, 24, 18, 9) + (2, -1, -2, 4, 9) \\ &= (33, -8, 30, 18, 6) 
\end{align*}

b/ Chọn $c_1 = 4, c_2 = -2, c_3 = -5$, ta có
một tổ hợp tuyến tính có từ \emph{S} là:

$\alpha = c_1\alpha_1 + c_2\alpha_2 + c_3\alpha_3$

\begin{align*}
 &= 4(-5, 2, -4, 2, 6) - 2(7, -1, 8, 6, 3) - 5(2, -1, -2, 4, 9) \\ &= \dots \\ &= \dots 
\end{align*}

\ul{Ví dụ mẫu 1}: bài 3.1.

\textbf{Bài 3.1.} Trong không gian tuyến tính $\mathbb{R}^3$ cho hệ $\{a_1, a_2, a_3\}$ với \[ a_1 = (1, 1, -1), \quad a_2 = (3, 2, 1), \quad a_3 = (-1, 1, 3). \] Chứng minh rằng phần tử $x = (7, 7, 3)$ là một tổ hợp tuyến tính của hệ $\{a_1, a_2, a_3\}$.

Chứng minh rằng
$x = (7, 7, 3)$ là một THTT
của $\{a_1, a_2, a_3\}$.

\ul{Giải}:

Để chứng tỏ $x = (7, 7, 3)$ là
một THTT của $\{a_1, a_2, a_3\}$
thì ta cần tìm $c_1, c_2, c_3 \in \mathbb{R}$
sao cho

$x = c_1a_1 + c_2a_2 + c_3a_3$

$\Leftrightarrow (7, 7, 3) = c_1(1, 1, -1) + c_2(3, 2, 1) + c_3(-1, 1, 3)$

\[ \Leftrightarrow (7, 7, 3) = (c_1, c_1, -c_1) + (3c_2, 2c_2, c_2) + (-c_3, c_3, 3c_3) \] \\ \[ \Leftrightarrow \begin{cases} c_1 + 3c_2 - c_3 = 7 \\ c_1 + 2c_2 + c_3 = 7 \\ -c_1 + c_2 + 3c_3 = 3 \end{cases} \]

Viết lại hệ pt dưới dạng ma trận hóa, ta có:

\[ \begin{pmatrix} 1 & 3 & -1 & | & 7 \\ 1 & 2 & 1 & | & 7 \\ -1 & 1 & 3 & | & 3 \end{pmatrix} \xrightarrow[(3) \rightarrow (3) + (1)]{(2) \rightarrow (2) - (1)} \begin{pmatrix} 1 & 3 & -1 & | & 7 \\ 0 & -1 & 2 & | & 0 \\ 0 & 4 & 2 & | & 10 \end{pmatrix} \xrightarrow{(3) \rightarrow (3) + 4(2)} \begin{pmatrix} 1 & 3 & -1 & | & 7 \\ 0 & -1 & 2 & | & 0 \\ 0 & 0 & 10 & | & 10 \end{pmatrix} \]

Do ta thu được 3 cột bán chuẩn liên tiếp nhau nên hệ pt có nghiệm duy
nhất là:

\[ \begin{cases} c_1 + 3c_2 - c_3 = 7 \\ -c_2 + 2c_3 = 0 \\ 10c_3 = 10 \end{cases} \Rightarrow \begin{cases} c_1 = 7 - 3c_2 + c_3 = 7 - 3.2 + 1 = 2 \\ c_2 = 2c_3 = 2 \\ c_3 = 1 \end{cases} \]

Vậy hệ luôn có nghiệm là
$c_1 = 2, c_2 = 2, c_3 = 1$ nên
$x = (7, 7, 3)$ là một THTT
của $\{a_1, a_2, a_3\}$.

\ul{Ví dụ mẫu 2}: bài 3.4

\textbf{Bài 3.4.} Trong không gian tuyến tính $\mathbb{R}^3$ cho hệ $\{a_1, a_2, a_3, a_4\}$ với $a_1 = (1, 1, 2), a_2 = (2, 3, -1), a_3 = (3, -1, 2), a_4 = (2, 8, -2)$. Hãy tìm tất cả các biểu diễn tuyến tính có thể có của $a_4$ trên hệ $\{a_1, a_2, a_3, a_4\}$.

Hãy tìm tất cả các biểu diễn tuyến tính có thể có của
$\alpha_4$ trên hệ
$\{a_1, a_2, a_3, a_4\}$.

\ul{Giải}:

Để tìm biểu diễn tuyến tính của
$\alpha_4$ trên hệ
$\{a_1, a_2, a_3, a_4\}$, nghĩa là ta
cần tìm $c_1, c_2, c_3, c_4 \in \mathbb{R}$ sao
cho

\[ a_4 = c_1a_1 + c_2a_2 + c_3a_3 + c_4a_4 \]

\[ \Leftrightarrow (2, 8, -2) = c_1(1, 1, 2) + c_2(2, 3, -1) + c_3(3, -1, 2) + c_4(2, 8, -2) \]

\[ \Leftrightarrow \begin{cases} c_1 + 2c_2 + 3c_3 + 2c_4 = 2 \\ c_1 + 3c_2 - c_3 + 8c_4 = 8 \\ 2c_1 - c_2 + 2c_3 - 2c_4 = -2 \end{cases} \]

Ta có viết dưới dạng ma trận hóa như sau:

\[ \begin{pmatrix} 1 & 2 & 3 & 2 & | & 2 \\ 1 & 3 & -1 & 8 & | & 8 \\ 2 & -1 & 2 & -2 & | & -2 \end{pmatrix} \xrightarrow[(3) \rightarrow (3) - 2(1)]{(2) \rightarrow (2) - (1)} \begin{pmatrix} 1 & 2 & 3 & 2 & | & 2 \\ 0 & 1 & -4 & 6 & | & 6 \\ 0 & -5 & -4 & -6 & | & -6 \end{pmatrix} \]

\[ \xrightarrow{(3) \rightarrow (3) + 5(2)} \begin{pmatrix} 1 & 2 & 3 & 2 & | & 2 \\ 0 & 1 & -4 & 6 & | & 6 \\ 0 & 0 & -24 & 24 & | & 24 \end{pmatrix} \xrightarrow{(3) \rightarrow -\frac{1}{24}(3)} \begin{pmatrix} 1 & 2 & 3 & 2 & | & 2 \\ 0 & 1 & -4 & 6 & | & 6 \\ 0 & 0 & 1 & -1 & | & -1 \end{pmatrix} \]

Do cột 4 không bán chuẩn hóa được nên hệ có vô số nghiệm như sau:

Đặt $c_4 = a, \forall a \in \mathbb{R}$

Ta có:

\[ \begin{cases} c_1 + 2c_2 + 3c_3 + 2c_4 = 2 \\ c_2 - 4c_3 + 6c_4 = 6 \\ c_3 - c_4 = -1 \end{cases} \Rightarrow \begin{cases} c_1 = 2 - 2c_2 - 3c_3 - 2c_4 = 2 - 2(2 - 2a) - 3(a - 1) - 2a = 1 - a \\ c_2 = 6 + 4c_3 - 6c_4 = 6 + 4(a - 1) - 6a = 2 - 2a \\ c_3 = c_4 - 1 = a - 1 \end{cases} \]

Như vậy, ta có thể viết lại

\[ a_4 = c_1a_1 + c_2a_2 + c_3a_3 + c_4a_4 \]

\[ = (1 - a)a_1 + (2 - 2a)a_2 + (a - 1)a_3 + a a_4 \]

Đây là tập hợp tất cả các biểu diễn tuyến tính của
$\alpha_4$ trên hệ
$\{a_1, a_2, a_3, a_4\}$, với mọi
$a \in \mathbb{R}$.

\ul{Ví dụ mẫu 3}:

Tìm $\lambda$ để
$x = (2, 3, 2, \lambda)$ biểu diễn được
theo các véc tơ

\[ a_1 = (1, 1, 2, 2), a_2 = (2, 3, 1, 4), a_3 = (3, 4, 2, 3) \]

\ul{Giải}:

Để $x = (2, 3, 2, \lambda)$ biểu diễn
được theo các véc tơ
$\alpha_1, \alpha_2, \alpha_3$ nghĩa là ta
cần tìm $c_1, c_2, c_3 \in \mathbb{R}$ sao
cho:

$x = c_1\alpha_1 + c_2\alpha_2 + c_3\alpha_3$ có nghiệm

$(2, 3, 2, \lambda) = c_1(1, 1, 2, 2) + c_2(2, 3, 1, 4) + c_3(3, 4, 2, 3)$ có nghiệm

\[ \Leftrightarrow \begin{cases} c_1 + 2c_2 + 3c_3 = 2 \\ c_1 + 3c_2 + 4c_3 = 3 \\ 2c_1 + c_2 + 2c_3 = 2 \\ 2c_1 + 4c_2 + 3c_3 = \lambda \end{cases} \]

có nghiệm

Ta viết lại hệ dưới dạng ma trận hóa như sau:

\[ \begin{pmatrix} 1 & 2 & 3 & | & 2 \\ 1 & 3 & 4 & | & 3 \\ 2 & 1 & 2 & | & 2 \\ 2 & 4 & 3 & | & \lambda \end{pmatrix} \xrightarrow[(4) \rightarrow (4) - 2(1)]{(3) \rightarrow (3) - 2(1)}{(2) \rightarrow (2) - (1)} \begin{pmatrix} 1 & 2 & 3 & | & 2 \\ 0 & 1 & 1 & | & 1 \\ 0 & -3 & -4 & | & -2 \\ 0 & 0 & -3 & | & \lambda - 4 \end{pmatrix} \xrightarrow{(3) \rightarrow (3) + 3(2)} \begin{pmatrix} 1 & 2 & 3 & | & 2 \\ 0 & 1 & 1 & | & 1 \\ 0 & 0 & -1 & | & 1 \\ 0 & 0 & -3 & | & \lambda - 4 \end{pmatrix} \] \\ \[ \xrightarrow{(4) \rightarrow (4) - 3(3)} \begin{pmatrix} 1 & 2 & 3 & | & 2 \\ 0 & 1 & 1 & | & 1 \\ 0 & 0 & -1 & | & 1 \\ 0 & 0 & 0 & | & \lambda - 7 \end{pmatrix} \]

Để hệ pt có nghiệm thì pt(4) phải thỏa dạng
$\lambda - 7 = 0 \Leftrightarrow \lambda = 7$ nghĩa là ta
có: $\lambda = 7$

Đáp số: $\lambda = 7$ là giá
trị cần tìm.

\textbf{4/ \ul{SỰ ĐỘC LẬP TUYẾN TÍNH (LINEARLY INDEPENDENCE) VÀ SỰ PHỤ
THUỘC TUYẾN TÍNH (LINEARLY DEPENDENCE)}}

Cho $(V, +, \bullet)$ là không
gian tuyến tính tổng quát trên trường số \emph{F},

và cho tập hợp $S = \{\alpha_1, \alpha_2, \dots, \alpha_m\}$
trên \emph{V}.

Ta nói \emph{S} là tập hợp độc lập tuyến tính (ĐLTT), nghĩa là các véc
tơ trong \emph{S} không có mối liên hệ gì với nhau, không phụ thuộc
nhau, và một véc tơ trong \emph{S} không thể biểu diễn theo các véc tơ
còn lại trong \emph{S}, nếu hệ pt sau:

$c_1\alpha_1 + c_2\alpha_2 + \dots + c_m\alpha_m = O$

(với $c_1, c_2, \dots, c_m$ là các ẩn
số cần tìm)

chỉ có bộ nghiệm tầm thường duy nhất là
$c_1 = c_2 = \dots = c_m = 0$.

Ngược lại, nếu hệ pt này có vô số nghiệm (bên cạnh bộ nghiệm tầm thường
$c_1 = c_2 = \dots = c_m = 0$) thì ta nói
tập hợp \emph{S} là phụ thuộc tuyến tính (PTTT), nghĩa là các véc tơ
trong \emph{S} có thể biểu diễn lẫn nhau, ví dụ
$\alpha_3 = 5\alpha_2 - 4\alpha_1 + 3\alpha_4$ hay
$\alpha_2 = 2\alpha_1 - \alpha_4 + 5\alpha_6$.

\ul{Ví dụ mẫu 4}:

Trên $V = \mathbb{R}^3$, cho tập
hợp \[ S = \{\alpha_1 = (1; -2; 1), \alpha_2 = (-1; 3; 2), \alpha_3 = (2; 0; 1), \alpha_4 = (-2; 1; 2)\} \].

Hỏi \emph{S} là ĐLTT hay PTTT, vì sao?

\ul{Giải}:

Để kiểm tra xem \emph{S} là ĐLTT hay PTTT, ta giải hệ pt

$c_1\alpha_1 + c_2\alpha_2 + c_3\alpha_3 + c_4\alpha_4 = O$

\begin{quote}
(ở đây ẩn số là
$c_1, c_2, c_3, c_4$)
\end{quote}

\[ \Leftrightarrow c_1(1, -2, 1) + c_2(-1, 3, 2) + c_3(2, 0, 1) + c_4(-2, 1, 2) = (0, 0, 0) \]

Viết lại hệ pt dưới dạng ma trận hóa ta có:

\[ \begin{pmatrix} 1 & -1 & 2 & -2 & | & 0 \\ -2 & 3 & 0 & 1 & | & 0 \\ 1 & 2 & 1 & 2 & | & 0 \end{pmatrix} \xrightarrow[(3) \rightarrow (3) - (1)]{(2) \rightarrow (2) + 2(1)} \begin{pmatrix} 1 & -1 & 2 & -2 & | & 0 \\ 0 & 1 & 4 & -3 & | & 0 \\ 0 & 3 & -1 & 4 & | & 0 \end{pmatrix} \xrightarrow{(3) \rightarrow (3) - 3(2)} \begin{pmatrix} 1 & -1 & 2 & -2 & | & 0 \\ 0 & 1 & 4 & -3 & | & 0 \\ 0 & 0 & -13 & 13 & | & 0 \end{pmatrix} \]

Do cột 4 không bán chuẩn hóa được nên hệ pt có vô số nghiệm như sau:

Đặt $c_4 = t, \forall t \in \mathbb{R}$

Ta có hệ

\begin{quote}
\[ \begin{cases} c_1 - c_2 + 2c_3 - 2c_4 = 0 \\ c_2 + 4c_3 - 3c_4 = 0 \\ -13c_3 + 13c_4 = 0 \end{cases} \Rightarrow \begin{cases} c_1 = c_2 - 2c_3 + 2c_4 = -t \\ c_2 = 3c_4 - 4c_3 = -t \\ c_3 = c_4 = t \end{cases} \]
\end{quote}

Cho nên ta nói \emph{S} là tập hợp PTTT.

\ul{Ví dụ mẫu 5}: Trên
$V = \mathbb{R}^4$, cho tập hợp
\[ S = \{\alpha_1 = (1; 0; 1; 0), \alpha_2 = (-1; 1; 0; 1), \alpha_3 = (0; 2; 1; -3)\} \].

Hỏi \emph{S} là tập hợp ĐLTT hay PTTT, vì sao?

\ul{Giải}:

Để kiểm tra \emph{S} là tập hợp ĐLTT hay PTTT, ta giải hệ pt

\[ c_1\alpha_1 + c_2\alpha_2 + c_3\alpha_3 = O \]

\begin{quote}
(ở đây ẩn số là
$C_1, C_2, C_3$)
\end{quote}

\[ \Leftrightarrow c_1(1, 0, 1, 0) + c_2(-1, 1, 0, 1) + c_3(0, 2, 1, -3) = (0, 0, 0, 0) \] \\ \[ \Leftrightarrow \begin{cases} c_1 - c_2 = 0 \\ c_2 + 2c_3 = 0 \\ c_1 + c_3 = 0 \\ c_2 - 3c_3 = 0 \end{cases} \]

Ta viết lại hệ pt dưới dạng ma trận hóa như sau:

\[ \begin{pmatrix} 1 & -1 & 0 & | & 0 \\ 0 & 1 & 2 & | & 0 \\ 1 & 0 & 1 & | & 0 \\ 0 & 1 & -3 & | & 0 \end{pmatrix} \xrightarrow{(3) \rightarrow (3) - (1)} \begin{pmatrix} 1 & -1 & 0 & | & 0 \\ 0 & 1 & 2 & | & 0 \\ 0 & 1 & 1 & | & 0 \\ 0 & 1 & -3 & | & 0 \end{pmatrix} \xrightarrow[(4) \rightarrow (4) - (2)]{(3) \rightarrow (3) - (2)} \begin{pmatrix} 1 & -1 & 0 & | & 0 \\ 0 & 1 & 2 & | & 0 \\ 0 & 0 & -1 & | & 0 \\ 0 & 0 & -5 & | & 0 \end{pmatrix} \]

Do hệ sau cùng có 3 cột bán chuẩn liên tiếp nhau, nên ta kết luận hệ có
nghiệm duy nhất

\[ \xrightarrow{(4) \rightarrow (4) - 5(3)} \begin{pmatrix} 1 & -1 & 0 & | & 0 \\ 0 & 1 & 2 & | & 0 \\ 0 & 0 & -1 & | & 0 \\ 0 & 0 & 0 & | & 0 \end{pmatrix} \] \\ \[ \begin{cases} c_1 - c_2 = 0 \\ c_2 + 2c_3 = 0 \\ -c_3 = 0 \end{cases} \Rightarrow c_1 = c_2 = c_3 = 0 \]

(là nghiệm tầm thường).

Nên ta nói \emph{S} là tập hợp ĐLTT.

* \emph{\textbf{\ul{CÁCH NHẬN DIỆN NHANH TẬP HỢP S LÀ ĐLTT HAY PTTT}}}:

Cho $(V, +, \bullet)$ là không
gian tuyến tính tổng quát trên trường số \emph{F},

và cho tập hợp $S = \{\alpha_1, \alpha_2, \dots, \alpha_m\}$
trên \emph{V}.

Để kiểm tra xem \emph{S} là ĐLTT hay PTTT, ta làm như sau:

+ Lập ma trận \[ A_S = \begin{pmatrix} \alpha_1 \\ \alpha_2 \\ \vdots \\ \alpha_m \end{pmatrix} \]

dòng 1

dòng 2

dòng \emph{m}

+ \ul{Trường hợp 1}:

$A_S$ tạo thành một
ma trận vuông cấp \emph{m}

Nếu $\det(A_S) \neq 0$ thì ta nói
\emph{S} là tập hợp ĐLTT.

Nếu $\det(A_S) = 0$ thì ta nói
\emph{S} là tập hợp PTTT.

+ \ul{Trường hợp 2}:

$A_S$ tạo thành một
ma trận tùy ý (có thể là ma trận vuông hay không vuông đều được).

Từ $A_S$ bán chuẩn
hóa (chuẩn hóa) tối đa các cột
$H$

Ta xác định $r(A_S) = rank(A_S) = $số
dòng khác zero của \emph{H}.

Nếu $r(A_S) = m \Rightarrow $số lượng
véc tơ trong
\emph{S}$\Rightarrow S$ là tập
hợp ĐLTT.

Nếu $r(A_S) < m \Rightarrow $số lượng
véc tơ trong
\emph{S}$\Rightarrow S$ là tập
hợp PTTT.

\ul{Ví dụ mẫu 6}: bài 3.9a

Chứng minh rằng hệ
$S = \{\alpha_1, \alpha_2, \alpha_3\}$ là hệ ĐLTT.

\ul{Giải}:

Lập ma trận

\begin{quote}
\[ A_S = \begin{pmatrix} \alpha_1 \\ \alpha_2 \\ \alpha_3 \end{pmatrix} = \begin{pmatrix} -2 & 1 & 1 \\ 1 & -2 & 1 \\ 1 & 1 & 2 \end{pmatrix} \]
\end{quote}

\ul{Cách 1}: Ta có

\begin{quote}
\[ \det(A_S) = (-2).(-2).2 + 1.1.1 + 1.1.1 - [1.(-2).1 + 1.1.2 + (-2).1.1] = 12 \neq 0 \]
\end{quote}

Suy ra hệ $\{\alpha_1, \alpha_2, \alpha_3\}$ là
hệ ĐLTT.

\ul{Cách 2}: Ta có

\begin{quote}
\[ A_S = \begin{pmatrix} -2 & 1 & 1 \\ 1 & -2 & 1 \\ 1 & 1 & 2 \end{pmatrix} \xrightarrow[(3) \rightarrow 2(3) + (1)]{(2) \rightarrow (2) - (3)} \begin{pmatrix} -2 & 1 & 1 \\ 0 & -3 & -1 \\ 0 & 3 & 5 \end{pmatrix} \xrightarrow{(3) \rightarrow (3) + (2)} \begin{pmatrix} -2 & 1 & 1 \\ 0 & -3 & -1 \\ 0 & 0 & 4 \end{pmatrix} = H \]
\end{quote}

Do ma trận \emph{H} có 3 dòng khác zero nên
$r(A_S) = rank(A_S) = 3 = $số véc tơ trong
\emph{S}.

Suy ra hệ $\{\alpha_1, \alpha_2, \alpha_3\}$ là
hệ ĐLTT.

\ul{Ví dụ mẫu 7}: bài 3.11a

Tìm các giá trị của
$\lambda$ để hệ
$\{\alpha_1, \alpha_2, \alpha_3\}$ là hệ ĐLTT.

\ul{Giải}:

Ta lập ma trận

\begin{quote}
\[ A_S = \begin{pmatrix} a_1 \\ a_2 \\ a_3 \end{pmatrix} = \begin{pmatrix} 1 & -1 & -1 \\ 1 & 2 & 3 \\ 2 & 1 & \lambda \end{pmatrix} \xrightarrow[(3)' \rightarrow (3)' + (1)']{(2)' \rightarrow (2)' + (1)'} \begin{pmatrix} 1 & 0 & 0 \\ 1 & 3 & 4 \\ 2 & 3 & \lambda + 2 \end{pmatrix} \]
\end{quote}

\[ \Rightarrow \det(A_S) = 1.c_{11}^{A_S'} = 1.(-1)^{1+1}.\begin{vmatrix} 3 & 4 \\ 3 & \lambda + 2 \end{vmatrix} = 3(\lambda + 2) - 12 = 3\lambda + 6 - 12 = 3\lambda - 6 \]

Để $\{\alpha_1, \alpha_2, \alpha_3\}$ là hệ ĐLTT
thì ta có $\det(A_S) \neq 0 \Leftrightarrow 3\lambda - 6 \neq 0 \Leftrightarrow \lambda \neq 2$.

Đáp số: $\lambda \neq 2$ thì hệ
$\{\alpha_1, \alpha_2, \alpha_3\}$ là hệ ĐLTT.

\end{document}
